Bide+ is a bus arrival prediction method that uses only the open data of the operator: the GPS positions of its
vehicles and its timetable. It is documented in a published technical report and an open reference implementation. This
paper tests it against the strongest available baseline, the operator's own real-time information, on the interurban bus
network of Biscay (Bizkaibus, Spain). For one week in November 2026 we recorded the official waiting times of 60 stops
every minute and the live positions every 20 seconds, and compared three predictions for the same bus at the same stop and
minute (the published timetable, the official real-time announcement and Bide+) against the actual passage reconstructed
from the feed. Over 160,046 paired predictions on five normal days, Bide+ is within two minutes of the actual passage in
64\% of cases, against 36\% for the official real time and 16\% for the timetable. With the bus 6 to 10 minutes away,
71\% of Bide+ predictions are within two minutes against 37\% of the official ones, and 61\% against 27\% at 10 to 20
minutes; the 95\% confidence intervals of the difference, by trip, are 33 to 36 points. Bide+ is more accurate at 58 of the
60 stops, at every hour and on every day. The official real time mostly announces buses later than they actually come, which makes
passengers miss them, while Bide+ errors are centred. The data, the capture program and the analysis are openly available.
